The rapid advancement of artificial intelligence (AI) has transformed agricultural systems from input-driven production models into data-driven decision ecosystems. While smart agriculture technologies have primarily focused on optimizing physical processes such as crop monitoring, irrigation, and yield prediction, comparatively less attention has been devoted to the development of AI-driven systems that support human capital formation---a critical determinant of long-term agricultural productivity and resilience.

Agricultural skill development programs represent one of the most widely implemented instruments for improving farmer performance. However, their effectiveness remains constrained by the uniform allocation of training interventions across heterogeneous farmer populations. Farmers differ systematically in baseline competencies, risk perception, resource access, and learning readiness, yet standardized curricula are commonly applied across all participants. This mismatch results in inefficient resource allocation and suboptimal learning outcomes, highlighting the need for more adaptive and intelligent allocation mechanisms \cite{anderson2004,feder1985}.

Recent advances in machine learning (ML) \cite{kamilaris2018,liakos2018,yang2025} enable the modeling of complex and non-linear relationships within agricultural datasets, offering new opportunities to predict heterogeneous outcomes across individuals. However, predictive capability alone is insufficient for real-world deployment. Decision-making in agricultural extension requires transparency, interpretability, and actionable logic that can be operationalized by practitioners. This has led to increasing interest in explainable artificial intelligence (XAI)\cite{lundberg2017,arrieta2020,ryo2022}, particularly SHapley Additive exPlanations (SHAP), which provides theoretically grounded feature attribution for tree-based models \cite{lundberg2017}. Despite these advances, a critical gap remains: the lack of integrated AI frameworks that translate predictive and explainable insights into structured and deployable decision support systems for agricultural training design.

Artificial intelligence has been increasingly applied to automate decision-making processes in agriculture, including quality inspection, prediction, and classification tasks \cite{yang2025,wonggasem2024}.

Deep learning models have demonstrated superior performance compared to traditional approaches in agricultural image-based applications \cite{wonggasem2024}.

This study addresses this gap by proposing an explainable AI-based Decision Support System (DSS) for precision allocation of agricultural skill development programs. The framework integrates K-Means clustering for farmer segmentation, Random Forest modeling for skill outcome prediction, and SHAP analysis for interpretable attribution, organized within a unified Segment--Attribute--Recommend pipeline. Unlike prior studies that treat explainability as a post-hoc interpretive tool, this study operationalizes SHAP outputs as the foundation for generating explicit decision rules, enabling the transition from predictive analytics to actionable decision intelligence.

The study is guided by three research questions: (1) Can unsupervised learning methods identify structurally meaningful farmer segments for training design? (2) Can machine learning models provide sufficient predictive performance to support individual-level decision-making in skill development? and (3) how can SHAP-based explanations be transformed into interpretable and implementable decision rules for agricultural extension systems?

This study makes three primary contributions to the field of artificial intelligence in agriculture. First, it extends the role of AI beyond predictive analytics toward actionable decision-making by proposing an integrated explainable AI-based decision support framework for human capital development. Second, it advances explainable AI by introducing a cluster-level SHAP approach that links model interpretability with structured decision logic. Third, it proposes a deployable Segment--Attribute--Recommend architecture that transforms model outputs into transparent and auditable decision rules, enabling practical implementation by domain practitioners.

Importantly, this study does not aim to establish causal relationships between training interventions and outcomes. The dataset is observational and non-experimental in nature; therefore, all findings are interpreted as associative. The framework is intended to support data-driven decision-making under real-world constraints rather than causal inference.

The remainder of the paper is structured as follows. Section~2 reviews the relevant literature. Section~3 presents the data and methodology. Section~4 reports the empirical results. Section~5 discusses the implications for AI-driven agricultural systems. Section~6 concludes the study.